% ===========================================================================
% PART 0 -- the statistics. Every concept gets five things, because a definition
% alone tells a reader what to write down, not what the thing IS:
%   the problem it solves | the picture | the definition | one concrete number |
%   how it is misread
% ===========================================================================
\section{The statistics, and what each concept actually is}
\label{sec:math}

\why{A definition is a rule for writing a symbol. It does not say why the object
exists, what it feels like, or what breaks without it. Each concept below is
therefore given the problem it was invented to solve, a picture, the definition,
one worked number, and the misreading that is standard in this literature.}

% ---------------------------------------------------------------------------
\subsection{Probability}

\prob{We need to talk about ``how often'' before we have seen everything. Without
a rule, ``often'' is a word; with one, it is arithmetic that composes.}

\intu{A finite bag of outcomes with weights that sum to $1$. An event is a
handful of outcomes; its probability is the weight of the handful. Nothing is
mysterious or infinite --- every probability in this paper is a fraction of a
finite set.}

\begin{definition}[Probability space]
A finite set $\Omega$ of outcomes with $\Prob(\omega)\ge0$ summing to $1$;
$\Prob(E):=\sum_{\omega\in E}\Prob(\omega)$ for $E\subseteq\Omega$.
\end{definition}

\feel{$968$ prompts, $42$ of which have a binding veto: the probability that a
randomly drawn prompt has one is $42/968 = 0.0434$.}

\misread{``Probability'' as a property of a single case. It is a property of a
\emph{procedure for drawing}. ``This prompt has probability $0.043$'' is not a
sentence; ``a prompt drawn uniformly from these $968$'' is.}

% ---------------------------------------------------------------------------
\subsection{The indicator, which turns yes/no into arithmetic}

\prob{Most questions here are yes/no --- did the decision change, is the veto
binding. Yes/no cannot be averaged, added, or given an interval.}

\intu{Write $1$ for yes and $0$ for no. Now the \emph{average} of the answers is
exactly the \emph{fraction} of yeses, so every tool built for averages applies
unchanged to proportions. This one move is why the whole paper needs only one
statistical toolkit.}

\begin{equation}
\frac1n\sum_{i=1}^n \mathbf{1}\{P_i\} = \frac{\#\{i: P_i \text{ true}\}}{n}.
\label{eq:indavg}
\end{equation}

\feel{Flip / no-flip over $968$ prompts, coded $1/0$: the average is $0.051$, and
that number is simultaneously ``the mean of the indicator'' and ``$5.1\%$ of
prompts flipped''. Same object.}

% ---------------------------------------------------------------------------
\subsection{Random variable}

\prob{We want to attach numbers to outcomes and then reason about the numbers.}

\intu{A rule that reads a number off whatever happened. It is neither random nor
a variable --- the randomness is in which outcome occurs, and the rule is fixed.
``$X$ is random'' is shorthand for ``the input to $X$ is drawn''.}

\begin{definition}
$X:\Omega\to\R$, and $\Prob(X=x):=\Prob(\{\omega: X(\omega)=x\})$.
\end{definition}

% ---------------------------------------------------------------------------
\subsection{Expectation}

\prob{One number summarising where a quantity sits, that does not depend on how
many times you looked.}

\intu{Take the long-run average. If value $x$ shows up a fraction $\Prob(X=x)$ of
the time, then in $n$ draws it contributes $x\cdot n\Prob(X=x)$; divide by $n$ and
the $n$ cancels. What remains does not depend on $n$, so it is a property of the
distribution rather than of the experiment. \emph{That} is why it is worth a
name.}

\begin{definition}
$\E[X] := \sum_x x\,\Prob(X=x)$.
\end{definition}

\begin{lemma}[Linearity]
\label{lem:lin}
$\E[aX+bY]=a\E[X]+b\E[Y]$, with \textbf{no independence assumption}.
\end{lemma}
\begin{proof}
Sum over outcomes rather than values --- legitimate because grouping a finite sum
differently does not change it. Then
$\sum_\omega (aX(\omega)+bY(\omega))\Prob(\omega)
= a\sum_\omega X(\omega)\Prob(\omega)+b\sum_\omega Y(\omega)\Prob(\omega)$.
\end{proof}

\intu{Linearity is the only tool in the paper that never needs a condition. Every
later step where an assumption would be contentious is routed through it
deliberately.}

\begin{corollary}[The bridge]
\label{cor:indexp}
$\E[\mathbf{1}\{P\}]=\Prob(P)$: a measured proportion estimates a probability.
\end{corollary}

\misread{``Expectation'' as what you expect to see. A fair die has expectation
$3.5$ and never shows it. It is a centre of mass, not a prediction.}

% ---------------------------------------------------------------------------
\subsection{Variance and standard deviation}

\prob{Two datasets can have the same average and be completely different. We need
a number for \emph{spread}, and the obvious candidate fails.}

\intu{Try averaging the deviations from the mean: it is \emph{exactly zero} for
every dataset, because positives cancel negatives. So the sign must be removed.
Squaring does it, and squaring has one property absolute value lacks: squared
spreads \emph{add} when quantities are independent, which is what makes the
spread of an average computable at all (\cref{eq:varmean}).}

\begin{definition}
$\Var(X):=\E[(X-\mu)^2]$; $\operatorname{sd}(X):=\sqrt{\Var(X)}$.
\end{definition}

\intu{The standard deviation is in the same units as the data --- it is a
\emph{typical distance from the middle}. The variance is in squared units and is
an intermediate quantity, which is why every interval in this paper is written
with the standard deviation and never the variance.}

\begin{lemma}[Computational form]
\label{lem:varalt}
$\Var(X)=\E[X^2]-(\E[X])^2$.
\end{lemma}
\begin{proof}
$\E[(X-\mu)^2]=\E[X^2-2\mu X+\mu^2]=\E[X^2]-2\mu^2+\mu^2$ by \cref{lem:lin}.
\end{proof}

\begin{corollary}[Variance of a yes/no quantity]
\label{cor:indvar}
For $X=\mathbf{1}\{P\}$ with $\Prob(P)=\pi$: since $0^2=0$ and $1^2=1$ we have
$X^2=X$, so $\Var(X)=\pi-\pi^2=\pi(1-\pi)$.
\end{corollary}

\feel{$\pi=0.5 \Rightarrow \Var=0.25$; $\pi=0.05 \Rightarrow \Var=0.0475$. A rare
event is intrinsically \emph{easier} to pin down than a coin flip, which is why
the $5\%$ flip rates in \cref{tab:power} have narrow intervals at modest $n$.}

\begin{lemma}[Scaling]
\label{lem:varscale}
$\Var(aX+b)=a^2\Var(X)$: shifting everything does not change spread; stretching
by $a$ stretches spread by $|a|$.
\end{lemma}
\begin{proof}
The deviation of $aX+b$ from its mean $a\mu+b$ is $a(X-\mu)$; square and take
expectations.
\end{proof}

\ex{Two prompts, both with mean satisfaction $0.5$ across the four responses.
Prompt~A: $(0.50,0.50,0.50,0.50)$ --- variance $0$, and \emph{no rubric can decide
anything here}, because every response is equally satisfying. Prompt~B:
$(0.95,0.55,0.45,0.05)$ --- variance $0.101$, and the decision is easy. Same mean,
opposite decidability. \textbf{The variance of the satisfaction vector is what
makes a criterion useful at all}, which is why a criterion that is constant across
responses contributes exactly nothing (measured at $0.0\%$ in \cref{tab:power}).}

\ex{Second use, at a different level. Across $968$ prompts the flip indicator has
$\pi=0.051$, so by \cref{cor:indvar} its variance is $0.0484$ --- and that single
number is what the entire interval $[4.4\%,5.8\%]$ is computed from. Nothing else
about the data enters.}

\misread{A large variance as ``bad data''. It is a fact about the quantity, not a
defect. What is bad is a large variance \emph{of the estimate}, which is a
different object (§\ref{sec:se}).}

% ---------------------------------------------------------------------------
\subsection{Independence, covariance, correlation}

\prob{We want to say ``these two move together'' without saying how strongly in
arbitrary units.}

\begin{definition}[Independence]
$\Prob(X=x,Y=y)=\Prob(X=x)\Prob(Y=y)$ for all $x,y$: knowing one tells you
nothing about the other.
\end{definition}

\intu{\textbf{Covariance:} multiply the two deviations and average. When both are
above their means the product is positive; when they move oppositely it is
negative; the average is therefore positive for co-movement. Its size is in the
product of the two units --- ``weight-points times satisfaction'' --- which is
uninterpretable, and that is the whole reason correlation exists.}

\begin{definition}
$\Cov(X,Y):=\E[(X-\mu_X)(Y-\mu_Y)]$; note $\Cov(X,X)=\Var(X)$.
\end{definition}

\begin{lemma}[Variance of a sum]
\label{lem:varsum}
$\Var(X+Y)=\Var(X)+\Var(Y)+2\Cov(X,Y)$; under independence the cross term is $0$.
\end{lemma}
\begin{proof}
$(D_X+D_Y)^2=D_X^2+2D_XD_Y+D_Y^2$ with $D_X:=X-\mu_X$; take expectations. Under
independence $\E[D_XD_Y]=\E[D_X]\E[D_Y]=0$.
\end{proof}

\begin{theorem}[Cauchy--Schwarz]
\label{thm:cs}
$|\Cov(X,Y)|\le \operatorname{sd}(X)\operatorname{sd}(Y)$.
\end{theorem}
\begin{proof}
For any $t$, $0\le\E[(D_X+tD_Y)^2]=\Var(X)+2t\Cov(X,Y)+t^2\Var(Y)$. A quadratic
that is never negative has discriminant $\le0$:
$4\Cov^2-4\Var(X)\Var(Y)\le0$.
\end{proof}

\begin{definition}[Correlation]
$r:=\Cov(X,Y)/(\operatorname{sd}(X)\operatorname{sd}(Y))\in[-1,1]$ by
\cref{thm:cs}.
\end{definition}

\intu{Correlation is covariance with the units divided out. $r$ answers: if one
quantity is one standard deviation above its mean, how many standard deviations
above is the other, on average? $r=0.3$ means ``about three tenths of one''.}

\feel{$r=+0.304$ between Shapley influence and alignment (§\ref{sec:c1}). Read:
a person one standard deviation more aligned than average is, on average, $0.30$
standard deviations more influential. That is a real relationship and a weak one,
which is exactly the finding.}

\ex{Concretely, in the three-person example of §\ref{sec:toy}: annotator~2 has
alignment $0.9313$ --- nearly perfect agreement with the collective outcome --- and
Shapley influence $0.3333$, less than half of annotator~1's $0.8333$ at a
\emph{lower} alignment of $0.7858$. Two people, and the two measures already rank
them differently. Across $1{,}865$ real pairs that shows up as $r=+0.304$:
the two quantities move together, weakly, and one cannot be substituted for the
other.}

\ex{What a \emph{high} correlation would have meant. Had $r$ come out $\approx1.0$,
the pre-registered conclusion was that influence and alignment are the same thing
and sixty rounds of correlational measurement had been fine. The number decided
between two research programmes, which is why it was registered before running.}

\misread{Three, all common.
(i) $r$ as a percentage --- it is not; only $r^2$ has a share reading, and only
under a linear model.
(ii) $r=0$ as ``unrelated'' --- covariance sees only \emph{linear} co-movement; a
quantity can be perfectly determined by another with $r=0$.
(iii) $r$ as comparable across studies --- it is scaled by the spreads present in
\emph{that} sample, so a restricted range shrinks it mechanically.}

% ---------------------------------------------------------------------------
\subsection{Estimator, estimand, bias}

\prob{We never see $\Prob$. We see $n$ numbers and must guess.}

\intu{The \emph{estimand} is the fixed unknown we want. The \emph{estimator} is a
recipe applied to data. The estimator wobbles from study to study; the estimand
does not move at all. Almost every misreading in applied statistics is a
confusion about which of the two is doing the moving.}

\begin{definition}
$\operatorname{Bias}(\hat\theta):=\E[\hat\theta]-\theta$.
\end{definition}

\intu{Unbiased means: right \emph{on average across repetitions of the whole
study}. It says nothing about being close in any one study. A recipe can be
unbiased and wildly variable, or biased and reliably close.}

\misread{``Unbiased'' as ``accurate''. They are different axes: bias is where the
centre sits, variance is how far it scatters.}

% ---------------------------------------------------------------------------
\subsection{The standard error}
\label{sec:se}

\prob{Our estimate came from one sample. How far might it be from the truth,
purely because a different sample would have given a different number?}

\intu{Averaging cancels noise. If each observation is off by a typical amount
$\sigma$, and the errors are independent, they partly cancel --- but only partly,
and the amount of cancellation is $\sqrt n$, not $n$, because errors add in
\emph{squares} (\cref{lem:varsum}) and lengths are square roots of squares. The
standard error is the typical distance between your estimate and the truth.}

\begin{align}
\E[\bar X] &= \mu &&\text{by \cref{lem:lin}}\notag\\
\Var(\bar X) &= \frac{1}{n^2}\sum_i\Var(X_i)=\frac{\sigma^2}{n}
  &&\text{by \cref{lem:varscale}, \cref{lem:varsum}}
\label{eq:varmean}\\
\se(\bar X) &= \sigma/\sqrt n \notag
\end{align}

\feel{$\sigma=0.124$, $n=968$: $\se=0.124/31.1=0.0040$. So a flip rate measured at
$5.1\%$ would typically land within about $0.4$ points of the truth --- and to
halve that to $0.2$ points you need \emph{four times} the prompts, not twice.}

\intu{The $\sqrt n$ is the entire economics of measurement. It is why
§\ref{sec:est} can say ``$1{,}213$ prompts for $1$pp resolution'' and mean it as a
budget, not a wish.}

\ex{The same measurement at three sample sizes, holding $\sigma=0.124$ fixed.
$n=10$: $\se=0.039$, interval $\pm7.7$ points --- a $5.1\%$ rate is
indistinguishable from $0\%$ or from $12\%$, so the study says nothing.
$n=100$: $\se=0.0124$, $\pm2.4$ points --- now $5\%$ is distinguishable from $0$
but not from $8\%$.
$n=968$: $\se=0.0040$, $\pm0.8$ points --- the actual study. \textbf{Same
phenomenon, same estimate, three different claims}, and only the third supports
the sentence ``deleting flips more often than doubling'' (a $0.93$ point gap).}

\ex{And the reason $\sqrt n$ hurts. To resolve $0.5$ points instead of $1$ point,
\cref{cor:n} demands $n=(2.8016\times0.124/0.005)^2=4{,}827$ prompts --- five times
the release, for twice the resolution. This is why §\ref{sec:design} states a
sample-size requirement rather than a wish.}

\misread{The standard error as the spread of the \emph{data}. It is the spread of
the \emph{estimate}, which is $\sqrt n$ times smaller. Data spread does not shrink
with more data; estimate spread does.}

\begin{remark}[Independence is doing real work here]
If observations are positively correlated, \cref{lem:varsum} adds covariance
terms and the true spread is \emph{larger} than $\sigma^2/n$. Using
\cref{eq:varmean} anyway understates uncertainty --- §\ref{sec:cluster} measures
exactly how much, because this error occurred in this work.
\end{remark}

% ---------------------------------------------------------------------------
\subsection{Why $n-1$}

\prob{To use $\se=\sigma/\sqrt n$ we need $\sigma$, which we also do not know.
Estimating it from the same data introduces a specific, correctable error.}

\intu{Deviations are measured from $\bar X$, which was computed \emph{from the
same points} and therefore sits closer to them than the true $\mu$ does. So the
squared deviations come out systematically too small. Dividing by $n-1$ instead
of $n$ inflates them by exactly the right amount.}

\begin{lemma}
\label{lem:bessel}
$\E\big[\sum_i (X_i-\bar X)^2\big]=(n-1)\sigma^2$.
\end{lemma}
\begin{proof}
$\sum_i(X_i-\bar X)^2=\sum_i(X_i-\mu)^2-n(\bar X-\mu)^2$ after expanding around
$\mu$ and using $\sum_i(X_i-\mu)=n(\bar X-\mu)$. Take expectations: the first term
is $n\sigma^2$, the second is $n\cdot\sigma^2/n=\sigma^2$ by \cref{eq:varmean}.
\end{proof}

\intu{``One degree of freedom was spent locating the centre.'' With $n=2$ points,
the mean sits exactly between them and the deviations carry only \emph{one}
piece of independent information about spread, not two --- hence $n-1=1$.}

% ---------------------------------------------------------------------------
\subsection{The central limit theorem}

\prob{To turn $\se$ into an interval we need to know the \emph{shape} of the
estimate's wobble, and the data here are $0/1$ --- about as far from a bell as
data gets.}

\intu{Averaging manufactures a bell. Whatever the individual observations look
like, their average is approximately normally distributed once $n$ is moderate,
because an average is a sum of many small independent contributions and sums of
many small independent things pile up in the middle. The data do not become
normal; the \emph{average} does.}

\begin{theorem}[Stated, not proved]
\label{thm:clt}
For iid $X_i$ with mean $\mu$ and finite variance,
$(\bar X-\mu)/(\sigma/\sqrt n)$ approaches the standard normal.
\end{theorem}

\begin{table}[H]\centering\small
\begin{tabularx}{\textwidth}{P{2.5cm}Y}
\toprule
licenses & normal-based intervals for \emph{averages}, including rates built from
$0/1$ indicators\\
does not say & that the data are normal --- they are not, and it does not matter\\
fails when & observations are dependent (§\ref{sec:cluster}); $n$ small and the
distribution skewed; or the quantity is not an average, e.g. a correlation --- which
is why §\ref{sec:boot} resamples instead\\
\bottomrule
\end{tabularx}
\end{table}

\begin{definition}[The three numbers used]
$z_q$ is the point with a fraction $q$ of the bell to its left:
$z_{0.975}=1.9600$, $z_{0.80}=0.8416$, $z_{0.99}=2.3263$. The bell is symmetric
about $0$, so $z_{1-q}=-z_q$.
\label{eq:quantiles}
\end{definition}

% ---------------------------------------------------------------------------
\subsection{Confidence intervals}
\label{sec:ci}

\prob{A single number with no indication of its precision is unusable: $5.1\%$
from ten prompts and $5.1\%$ from a thousand are different claims.}

\intu{Build a net around the estimate wide enough that, if you repeated the whole
study forever, $95\%$ of the nets would catch the fixed truth. \textbf{The net
moves; the truth does not.} Everything confusing about intervals comes from
imagining it the other way round.}

\begin{steps}
\item From \cref{thm:clt}, with probability $0.95$:
      $-1.96\le(\bar X-\mu)/\se\le1.96$.
\item Multiply through by $\se>0$; inequalities are preserved because the
      multiplier is positive: $-1.96\,\se\le \bar X-\mu\le 1.96\,\se$.
\item Subtract $\bar X$, then multiply by $-1$ --- which \textbf{reverses} both
      inequalities: $\bar X-1.96\,\se\;\le\;\mu\;\le\;\bar X+1.96\,\se$.
\end{steps}

\feel{$5.1\%$ with $\se=0.37$pp gives $[4.4\%,\,5.8\%]$. Read: a study like this
one, repeated, would produce intervals catching the truth $19$ times in $20$.}

\ex{Two studies reporting the identical estimate $5.1\%$.
Study~1: $n=25$, interval $[0.0\%,\,10.9\%]$ --- compatible with no effect at all.
Study~2: $n=968$, interval $[4.4\%,\,5.8\%]$ --- incompatible with anything below
$4$ points.
\textbf{The point estimate carries none of this.} Quoting $5.1\%$ without the
interval hides the entire difference between a result and a rumour.}

\misread{Three, in increasing order of consequence.
(i) ``There is a $95\%$ chance $\mu$ is in this interval'' --- $\mu$ is a constant;
it is in or it is not.
(ii) ``$95\%$ of future observations fall here'' --- that is a prediction interval
and is far wider.
(iii) ``The interval excludes zero, so the effect is large'' --- exclusion is
about \emph{precision}, not size. A trivial effect measured well excludes zero;
§\ref{sec:mde} separates the two questions.}

% ---------------------------------------------------------------------------
\subsection{Tests, $z$, and $p$}

\prob{We want a rule for ``this is more than nothing'' that cannot be tuned after
seeing the answer.}

\intu{Assume the boring world (the effect is zero). Ask how surprising the data
would be \emph{in that world}. If very surprising, the boring world is a poor
description. Note what this does \emph{not} do: it never evaluates the interesting
world, and it never assigns a probability to either.}

\begin{definition}
$z:=\hat\theta/\se$, the estimate measured in units of its own uncertainty;
$p:=\Prob(|Z|\ge|z|\mid H_0)$.
\end{definition}

\feel{$z=+3.3$ for the paired delete-versus-double difference. Read: the observed
gap is $3.3$ of its own standard errors away from zero.}

\misread{Three, all false.
(i) $p$ is not $\Prob(H_0\mid\text{data})$ --- it is conditioned the other way.
(ii) $p$ is not the probability the result is a fluke.
(iii) $1-p$ is not the probability the effect is real.
A $p$-value is a statement about \emph{data under an assumption}, never about the
assumption given data. This paper mostly reports $z$ and intervals, because a
$p$-value collapses size and precision into one number and they must stay
separate.}

% ---------------------------------------------------------------------------
\subsection{Power, and the minimum detectable effect}
\label{sec:mde}

\prob{A null result is uninterpretable on its own. ``We found nothing'' and ``we
could not have found anything'' produce the same output.}

\intu{Two bells: one centred at zero (the boring world), one centred at the true
effect. The test rejects when the estimate lands past a threshold in the first
bell's tail. \emph{Power} is how much of the \emph{second} bell sits past that
threshold. If the two bells overlap heavily --- a small effect, or a large
$\se$ --- most of the second bell is on the wrong side and the study will miss a
real effect most of the time. The MDE is the separation at which $80\%$ of the
second bell finally clears.}

\begin{theorem}
\label{thm:mde}
$\mathrm{MDE}=(z_{1-\alpha/2}+z_{1-\beta})\,\se$.
\end{theorem}
\begin{proof}
Reject when $\hat\theta\ge z_{1-\alpha/2}\se$. If the truth is $\delta$, then
$(\hat\theta-\delta)/\se$ is standard normal, so rejection requires
$(\hat\theta-\delta)/\se \ge z_{1-\alpha/2}-\delta/\se$. That has probability
$1-\beta$ exactly when the right side equals $-z_{1-\beta}$, by symmetry of the
bell. Solve for $\delta$.
\end{proof}

\feel{$\se=0.37$pp gives $\mathrm{MDE}=2.8016\times0.37=1.05$pp. Read: this design
would reliably notice a true flip rate difference of about one point and would
usually miss half a point.}

\begin{corollary}[Sample size]
\label{cor:n}
$n=(2.8016\,\sigma/\delta)^2$ to detect $\delta$.
\end{corollary}

\ex{A real case from this programme, and it explains a retraction. A round
claimed ``the compiled standard respects vetoes better than a human peer'', with a
measured gap of $1.72$ points; the claim was later withdrawn after it failed in
$7$ of $12$ held-out halves, and nobody asked why. Computing the MDE from the
round's own published interval: $\se=0.0069$, so
$\mathrm{MDE}=2.8016\times0.0069=1.93$ points. \textbf{The design's minimum
detectable effect was larger than the effect it was measuring.} The full-sample
interval barely excluding zero was luck; halving the data removed the luck. The
held-out sweep was rediscovering a power problem twelve times, and two lines of
arithmetic --- available from numbers already written down --- would have said so
at the start.}

\ex{The same computation used forwards. The paired delete-versus-double
difference is $+0.93$ points with $\mathrm{MDE}=0.79$. Effect exceeds MDE, but
only by a factor of $1.2$: the finding is real and its margin is thin, and that
is stated in \cref{tab:power} rather than buried.}

\misread{A null as evidence of absence. If $\mathrm{MDE}$ exceeds the effect the
null was meant to rule out, the null is \emph{silence}. Every null in
§\ref{sec:findings} carries its MDE for this reason.}

% ---------------------------------------------------------------------------
\subsection{Clustering, ICC, and the design effect}
\label{sec:cluster}

\prob{We have $19{,}360$ rows but only $968$ prompts, each perturbed $20$ times.
Treating the rows as independent claims twenty times more information than we
have. \textbf{This error was made in this work and had to be corrected.}}

\intu{Asking one person twenty questions is not the same as asking twenty people
one question. Repeated observations of the same prompt share whatever is peculiar
to that prompt, so the twentieth observation adds much less than the first. How
much less is exactly what the ICC measures.}

\begin{pts}
\item \textbf{The model.} $X_{pr}=\mu+a_p+e_{pr}$, with $\Var(a_p)=\sigma_b^2$
      (``between prompts'') and $\Var(e_{pr})=\sigma_w^2$ (``within a prompt'').
\item \textbf{The correct variance.} Cluster \emph{means} are independent, so
      \cref{eq:varmean} applies to them --- not to the rows.
\end{pts}
\begin{equation}
\Var(\hat\pi)=\frac1P\Big(\sigma_b^2+\frac{\sigma_w^2}{R}\Big).
\label{eq:varclust}
\end{equation}

\begin{definition}[Intraclass correlation]
$\rho:=\sigma_b^2/(\sigma_b^2+\sigma_w^2)$: the share of all variation that is
\emph{between} prompts. Equivalently, the correlation between two observations
drawn from the same prompt.
\end{definition}

\intu{$\rho=0$: repeats are as good as fresh prompts. $\rho=1$: all repeats of a
prompt are identical and the nineteen extra ones are worthless.}

\begin{theorem}[Design effect]
\label{thm:deff}
$\mathrm{DEFF}=\dfrac{\text{\cref{eq:varclust}}}{\sigma^2/(PR)}=1+(R-1)\rho$.
\end{theorem}
\begin{proof}
$\dfrac{\frac1P(\sigma_b^2+\sigma_w^2/R)}{(\sigma_b^2+\sigma_w^2)/(PR)}
=\dfrac{R\sigma_b^2+\sigma_w^2}{\sigma_b^2+\sigma_w^2}=1+(R-1)\rho$.
\end{proof}

\feel{Measured: $\sigma_b^2=0.01361$, $\sigma_w^2=0.03677$, so $\rho=0.2701$ and
with $R=20$, $\mathrm{DEFF}=6.13$. Read: the naive standard error is too small by
$\sqrt{6.13}=2.48$, every naive $z$ is inflated by the same factor, and
$19{,}360$ rows carry the information of about $3{,}150$. The honest count of
independent units stays $968$.}

\ex{Why the repeats are correlated here, mechanically. Each of the $20$ repeats
perturbs a \emph{different criterion of the same prompt}. If that prompt's four
responses are far apart in score --- a wide margin --- then \emph{no} single-criterion
perturbation flips it, and all twenty repeats return $0$ together. If the margin
is narrow, many return $1$ together. The prompt's margin is a property the
repeats \textbf{share}, and $\rho=0.27$ is the size of that sharing.}

\ex{What it cost. A pairwise-dependence screen reported ``$89$ of $120$ pairs
survive'' computed on the pooled $19{,}360$ rows. Recomputed at the cluster level:
$74$ of $120$. Fifteen apparent findings were the design effect. The derived
inflation is $\sqrt{6.13}=2.48$ and the directly resampled inflation was $2.79$;
they agree, and the residual is that a correlation is not a simple mean so
\cref{thm:deff} is only approximate for it.}

\misread{Row count as sample size. The question is always ``how many independent
things did I observe'', and repeated measurement of the same thing is not
independence.}

% ---------------------------------------------------------------------------
\subsection{The bootstrap}
\label{sec:boot}

\prob{The formulas above give $\se$ for an \emph{average}. For a correlation, a
median, a ratio, or a rank statistic there is generally no formula.}

\intu{Replace the algebra with imitation. We cannot rerun the study, but we can
rerun a \emph{stand-in} for it: draw a new dataset from the one we have, with
replacement, and recompute. The spread of the recomputed values imitates the
spread we would have seen across genuinely fresh studies, because the observed
sample is the best available stand-in for the population.}

\begin{protocol}[Cluster bootstrap]
\label{prot:boot}
Resample $P$ \emph{clusters} with replacement; recompute the statistic; repeat
$B$ times; report the $2.5$th and $97.5$th percentiles.
\end{protocol}

\intu{The resampling unit must be the cluster. Resampling rows would build
datasets whose internal dependence differs from the real one, reproducing exactly
the understatement of \cref{thm:deff}. Every interval in this paper resamples
prompts.}

\feel{$400$ resamples of $968$ prompts gave $[+0.2566,+0.3520]$ around the
observed $+0.3040$ in §\ref{sec:c1} --- a quantity for which no closed-form
$\se$ exists.}

\misread{Bootstrap intervals as assumption-free. They fail with few clusters, at
the edge of a parameter's range, and for statistics that are not smooth in the
data. All three are checked in §\ref{sec:limits}.}

% ---------------------------------------------------------------------------
\subsection{Permutation tests}

\prob{Sometimes the question is not ``how big'' but ``is the pairing between these
two lists doing any work at all''.}

\intu{Destroy the pairing and see if the result survives. Shuffle one list against
the other many times; if the real pairing gives something the shuffles rarely
give, the pairing carries information. It is a \emph{falsifier}: it can say the
pairing matters, never \emph{why}.}

\feel{$B=200$ shuffles: the smallest attainable $p$ is $1/(B+1)=0.005$. A reported
$p=0.0000$ from $200$ shuffles means $p<0.005$ and nothing more.}

\ex{What it answers here. Two operators each flip about $5\%$ of prompts. Do they
flip the \emph{same} prompts? Shuffling which prompt each operator's outcome is
attached to destroys any such coupling while leaving both flip \emph{rates}
untouched --- so the shuffled distribution is exactly the world where the two
operators are unrelated but equally active. Comparing the real coupling to that
distribution isolates the coupling and nothing else.}

\misread{The shuffle as neutral. Shuffling assumes all orderings are equally
likely under the null; if the list has internal structure --- clustering, time
order --- shuffling destroys that too, and the test answers a question nobody
asked. In this paper shuffles are applied to cluster-level values for exactly
this reason.}

% ---------------------------------------------------------------------------
\subsection{Multiplicity}

\prob{Test $170$ things at the usual threshold and about $8$ will look real when
nothing is.}

\intu{Every test is a lottery ticket for a false positive. Buy $170$ tickets and
you should expect to win. Either raise the bar per test, or change what you are
controlling.}

\begin{theorem}[Bonferroni]
\label{thm:bonf}
Testing $C$ hypotheses each at level $\alpha/C$ keeps the probability of
\emph{any} false rejection at $\le\alpha$.
\end{theorem}
\begin{proof}
$\mathbf{1}\{\bigcup A_i\}\le\sum_i\mathbf{1}\{A_i\}$ pointwise (\cref{lem:union});
take expectations via \cref{cor:indexp}: $\Prob(\bigcup A_i)\le C\cdot\alpha/C$.
\end{proof}

\feel{$C=171$, $\alpha=0.05$: per-test level $0.000292$, whose two-sided normal
quantile is $|z|=3.62$. This paper uses the stricter $|z|>3.9$ and always reports
tests-run beside tests-surviving.}

\begin{remark}[Bonferroni versus false discovery rate]
Bonferroni controls the chance of \emph{any} false positive --- right when one
false claim is costly. Benjamini--Hochberg instead controls the expected
\emph{share} of false positives among rejections, with threshold $qk/C$ at rank
$k$; at $k=C$ that is $q$ itself, not $q/C$. Confusing the two demands orders of
magnitude more resolution than BH requires.
\end{remark}

\ex{Why $170$ tests is not an unusual number. The operator grid in this work has
$19$ operators; asking which pairs behave alike is $19\times18/2=171$ questions,
and nobody set out to run $171$ tests --- they fell out of one table. At the usual
$5\%$ threshold about $8.5$ would look real with nothing there, which is within a
factor of two of the number of \emph{genuine} findings. The correction is not
pedantry; without it the table is half noise.}

\misread{Reporting only the survivors. That is the multiplicity failure with
manners: the correction is meaningless unless the denominator is published.}

% ---------------------------------------------------------------------------
\subsection{Measurement error, reliability, attenuation}
\label{sec:atten}

\prob{Every quantity here is read by a fallible instrument --- a language model
judging whether a criterion is satisfied. What does that do to the numbers?}

\intu{Measuring with a blurry ruler. Blur adds spread that has nothing to do with
the thing being measured, so the measured quantity varies more than the true one.
When you correlate two blurry measurements, the shared signal is unchanged but
both spreads are inflated --- and correlation is signal divided by spreads. So
\textbf{noise always pushes a correlation toward zero}, never away. It is a
systematic pull, not random error, and more data does not fix it.}

\begin{definition}[Reliability]
$\mathrm{rel}:=\Var(T)/\Var(X)=\sigma_T^2/(\sigma_T^2+\sigma_\eps^2)$: the share
of observed variation that is signal.
\end{definition}

\begin{lemma}[The identity that makes this free to compute]
\label{eq:relcorr}
For two independent measurements $X_1=T+\eps_1$, $X_2=T+\eps_2$ with equal error
variances, $\Corr(X_1,X_2)=\mathrm{rel}$.
\end{lemma}
\begin{proof}
$\Cov(X_1,X_2)=\Var(T)$, since every cross term involves an independent error;
each has variance $\sigma_T^2+\sigma_\eps^2$.
\end{proof}

\intu{\textbf{Two independent measurements of the same thing agree exactly as much
as one of them is trustworthy.} So running a second judge measures the first
judge's reliability at no extra conceptual cost --- which is what
§\ref{sec:instrument} does.}

\begin{theorem}[Attenuation]
\label{thm:atten}
$\Corr(X,Y)=\Corr(T,U)\sqrt{\mathrm{rel}_X\mathrm{rel}_Y}$.
\end{theorem}
\begin{proof}
$\Cov(X,Y)=\Cov(T,U)$; and $\operatorname{sd}(X)=\sigma_T/\sqrt{\mathrm{rel}_X}$
from the definition of reliability. Substitute.
\end{proof}

\feel{Measured judge reliability $0.5497$, so $\sqrt{0.5497}=0.741$. Read: a
\emph{perfect} true relationship, measured through this judge, reads $0.74$. No
sample size raises it --- the pull is bias, and $\sqrt n$ does not touch bias.}

\ex{The judge, concretely. Two language models scored the same $59{,}936$
criterion-by-response cells; they correlate at $0.5497$. By \cref{eq:relcorr} that
\emph{is} one judge's reliability, so by \cref{thm:atten} every correlation read
through this tensor is multiplied by $\sqrt{0.5497}=0.741$. A relationship that is
genuinely perfect reads $0.74$; one that genuinely reads $0.4$ was really $0.54$.
\textbf{Every number in this paper that routes through satisfaction inherits this
factor}, and the compounded ceiling with the panel's own reliability is $0.720$.}

\ex{And the use that matters. §\ref{sec:gaps} asks whether the two elicitation
routes disagree only because both are noisy. Dividing the observed $0.6265$ by
$\sqrt{0.8702\times0.9408}$ removes exactly the blur and leaves $0.6924$. The gap
does not close --- so the disagreement is between what people write and what they
choose, not between two noisy views of one thing. \textbf{Without
\cref{thm:atten} that question has no answer at all.}}

\misread{A weak measured correlation as a weak true relationship. It may be a
strong relationship seen through a blurry instrument, and \cref{thm:atten} is
the arithmetic that tells them apart --- used in §\ref{sec:gaps} to establish that
the two elicitations disagree \emph{beyond} their own blur.}

% ---------------------------------------------------------------------------
\subsection{Spearman--Brown: the reliability of a panel}

\prob{One judge is unreliable. A panel of many should be better. By how much,
exactly?}

\intu{Averaging $m$ measurements multiplies the signal by $m$ but the noise only
by $\sqrt m$, because independent noise adds in squares (\cref{lem:varsum}). So
signal-to-noise improves with panel size, but with diminishing returns --- and the
formula says precisely how fast.}

\begin{theorem}
\label{thm:sb}
Doubling the number of components takes reliability $r$ to $2r/(1+r)$.
\end{theorem}
\begin{proof}
Summing $m$ components gives signal variance $m^2\sigma_T^2$ and error variance
$m\sigma_\eps^2$, so $\mathrm{rel}_m=m\lambda/(m\lambda+1)$ with
$\lambda=\sigma_T^2/\sigma_\eps^2$. At $m=1$, $r=\lambda/(\lambda+1)$, so
$\lambda=r/(1-r)$; substitute $m=2$.
\end{proof}

\feel{Split-half $r=0.8961$ gives full-panel $2(0.8961)/1.8961=0.9452$. Read:
splitting the annotators of a prompt in half and correlating the two halves
measures a \emph{half}-sized panel; \cref{thm:sb} corrects it to the panel that
was actually used, which is the number that enters \cref{thm:atten}.}

% ---------------------------------------------------------------------------
\subsection{Two-way clustering}

\prob{Here observations are grouped two ways at once: an annotator appears on many
prompts, and a prompt is rated by many annotators. Correcting for one leaves the
other.}

\intu{Two observations are dependent if they share a prompt \emph{or} an
annotator. Add the two corrections and you have double-counted the pairs sharing
both, so subtract that once. It is inclusion--exclusion applied to covariance
terms.}

\begin{theorem}[Cameron--Gelbach--Miller]
\label{thm:cgm}
$\hat V_{\text{two-way}}=\hat V_{\text{prompt}}+\hat V_{\text{annotator}}
-\hat V_{\text{both}}$.
\end{theorem}

\misread{A slightly negative estimate as a bug. It signals that one dimension
carries almost no dependence, which is information.}
